\section{Knowledge Graphs for Code}
\label{sec:sota-kg}

Graph representations of programs have a long history as analysis artefacts, and a
short one as prompt material. The code property graph unified syntax, control flow
and data dependence into a queryable structure and demonstrated it on vulnerability
discovery~\cite{yamaguchi2014cpg}; learning directly over program graphs was shown
to beat learning over token sequences on tasks requiring reasoning about variable
use~\cite{allamanis2018programgraphs}. Both predate the current interest in
supplying graphs to language models, and both supply its justification: the
relationships that matter for correctness are often distant in the token stream and
adjacent in the graph.

Three recent systems build repository-scale graphs to serve a language model, and
each targets a different task. GraphCoder retrieves context for repository-level
code completion by traversing a code-context graph rather than by embedding
similarity, and reports gains over similarity retrieval on that
task~\cite{liu2024graphcoder}. KGCompass links repository artefacts, including
issues and pull requests, to code entities in one graph and uses paths through it to
localise faults for automated repair, reporting that the large majority of the
bugs it localises successfully are reachable only by traversing more than one edge
and are missed by approaches without the graph~\cite{yang2025kgcompass}. Prometheus
maintains a repository graph as persistent memory for an agent navigating a
long-horizon programming task~\cite{pan2025prometheus}. All three report their
results on repository-level benchmarks such as
SWE-bench~\cite{jimenez2024swebench}.

Two observations about this body of work shape what follows.

The first is the gap in task coverage. Completion, repair and navigation are all
tasks with a checkable outcome: the code compiles, the test passes, the patch
resolves the issue. Review-comment generation has no such outcome, which is
presumably why the graph literature has not addressed it. The same absence of
an automatic success signal also makes the problem worth addressing, since it is the
condition under which evaluation methodology becomes part of the contribution.

The second is that these systems' graphs are built by static analysis, so their
edges are resolved rather than guessed. This thesis follows the same approach: the
structural context in Chapters~\ref{chap:approach} and~\ref{chap:results} is
extracted from a code property graph produced by Joern, which resolves calls and
data-flow edges across files. The unexamined question is not what kind of graph to
build, but whether a resolved graph helps on a task---review-comment
generation---where no prior work has tested it.

\subsection{Change Impact Analysis}
\label{sec:sota-cia}

The dependency information this thesis supplies to the reviewer is, in
substance, a change impact analysis: given an edit, identify the programme
elements that may be affected~\cite{bohner1996cia}. The field has over
thirty years of work, from Weiser's programme
slicing~\cite{weiser1984slicing} through call-graph and data-flow
reachability~\cite{ryder2001cia} to modern test-impact analysis that
selects the subset of a test suite exercising changed
code~\cite{legunsen2016testimpact}. What is new in this thesis is not the
dependency computation---the CPG resolves the same edges these analyses
do---but the delivery mechanism: the impact set is rendered as natural
language and handed to a generative model rather than used as a programme
transformation. Experiment~2 isolates this delivery step by measuring how
much of the analyser's output the model preserves.

\subsection{Agentic Code Review}
\label{sec:sota-agentic}

Tool-using language models can navigate a repository by issuing file-open,
grep, and static-analysis commands during generation. Systems such as
SWE-agent~\cite{yang2024sweagent} and
Agentless~\cite{xia2024agentless} already do this for issue resolution.
Agentic code reviewers that browse the repository during review are a
natural extension.

This thesis takes the complementary approach: pre-computed structural
context supplied in a single prompt, with no tool calls during generation.
Pre-computation offers determinism (the context is frozen and reproducible),
bounded cost (one CPG build per PR, no iterative model calls), and
auditability (the evidence pack is a human-readable JSON file). The
trade-off is that the model cannot ask follow-up questions or explore paths
the builder did not anticipate. A direct comparison between pre-computed
context and agent-driven navigation is left to future work.
